\scp{Summary of previous studies in the region}{01-05}
The languages of this region have long been the subject of debate as to how they fit into the larger Austronesian picture.
The debate has tended to be marginal to the discussion of Austronesian as a whole,
and based on very little data because the languages in this region are severely under-documented.
Very few languages have published grammars or dictionaries,
although the past ten years has seen important additions.
There are topic-specific papers published on a few additional languages.
Most comparative work has been based on scattered word lists of varying length
and reliability, and often full of transcription errors
and glossing errors.\footnote{For example,
		Grimes counted 181 errors of transcription
		and glossing in an online 200+ item word list of the Buru language
		by another modern scholar not familiar with that language,
		even allowing for dialect differences and lexical taboos.
		\citet{Jonker1908,Jonker1915} used three different strategies for indicating double vowels in Rote languages,
		and this is often overlooked in the comparative literature,
		even though it is of comparative-historical significance.}
\citet{Clark1987} noted there may be 40--50 languages in the Central Malayo-Polynesian (CMP) group.
\citet[288]{Blust1993} estimates “between 90 and 100 CMP languages”.
\citet{Grimes1995} list 150 CMP languages by name with additional information on location,
number of speakers, etc.
\citet[32]{Blust2013} estimates CMP “includes about 120 languages”.
This present study addresses the classification of around 250 Austronesian languages.
These wildly varying figures show how little is known about the languages in this region.

The justification for talking about the Austronesian languages of eastern Indonesia
as some sort of subgroup(s) is based primarily on a series of independent observations
and assumptions going back at least as far as
\citet{Brandes1884,Friederici1913,Jonker1914,Esser1938} and \citet{Salzner1960},
with many referring to each other,
that the languages of Nusa Tenggara (the Lesser Sundas)
and Maluku are somehow different from those to the west
and north (western Indonesia,
Philippines)
and somehow related to those farther east (Melanesia).
However for the most part,
the arguments have been by declaration rather than by demonstration.
Of the scholars discussing these comparative issues,
only \citet{Jonker1908,Jonker1913,Jonker1915,Jonker1932}
and \citet{Stresemann1918,Stresemann1927} had a significant knowledge of the region in question,
having written grammars or dictionaries of languages in what is now eastern Indonesia.

Perhaps the simplest way to group scholars is on their view
of whether the Austronesian languages in present day eastern Indonesia
and Timor-Leste group with those to the west
and north or to the east,
or are transitional,
as summarised in \trf{01-01}.

\begin{table}
	\centering\tcp{Previous studies grouping languages with the west or the east}{01-01}
	\st{6pt}\begin{tabular}{ll}\lsptoprule
			\tbf{Group with the west}				& \tbf{Group with the east}	\\ \midrule
			\citet{Brandstetter1906}	& \citet{Brandes1884}	\\
			\citet{Dempwolff1934,Dempwolff1937,Dempwolff1938}	& \citet{Friederici1913}	\\
			\citet{Salzner1960}; Bima-Sumba										& \citet{Schmidt1926}	\\
			\citet{Haudricourt1965}		& \citet{Cappell1943,Cappell1944}	\\
			\citet{Dyen1965,Dyen1978a}& \citet{Salzner1960}; Maluku	\\
			\citet{Dahl1976}					& \citet{Blust1978,Blust1993,Blust2013}	\\ \midrule
			\mc{2}{c}{\tbf{Transitional}} \\ \midrule
			\mc{2}{c}{\citet{Kern1906}} \\
			\mc{2}{c}{\citet{Jonker1914}} \\
			\mc{2}{c}{\citet{KanskiKasprusch1931}} \\
		\lspbottomrule
	\end{tabular}
\end{table}

%\begin{table}
	%\centering\tcp{Previous studies grouping languages with the west or the east}{01-01}
	%\begin{tabular}{p{0.3\columnwidth}p{0.3\columnwidth}p{0.3\columnwidth}}\lsptoprule
%Group with the west	&	Transitional	&	Group with the east	\\ \midrule
%\citet{Brandstetter1906}	&	\citet{Kern1906}	&	\citet{Brandes1884}	\\
%\citeauthor{Dempwolff1934} (\citeyear{Dempwolff1934}, \citeyear{Dempwolff1937},	&	\citet{Jonker1914}	&	\citet{Friederici1913}	\\
%\citeyear{Dempwolff1938}) & \citeauthor{KanskiKasprusch1931} & \citet{Schmidt1926}	\\
%\citet{Salzner1960}; 	&	(\citeyear{KanskiKasprusch1931})	&	\citet{Cappell1943,Cappell1944}\\
%Bima-Sumba 	& & \citet{Salzner1960}; Maluku\\
%\citet{Haudricourt1965}	&		&		\citet{Blust1978,Blust1993,Blust2013}\\
%\citet{Dyen1965,Dyen1978a}	&		&		\\
%\citet{Dahl1976}	&		&		\\
		%\lspbottomrule
	%\end{tabular}
%\end{table}

%\begin{table}
	%\centering\caption{Previous studies grouping languages with the west or the east}
	%\label{tab:Previous studies grouping languages with the west or the east}
	%{\small{\wg\begin{tabular}{p{0.3\columnwidth}p{0.3\columnwidth}p{0.3\columnwidth}}\lsptoprule
%Group with the west	&	Transitional	&	Group with the east	\\ \midrule
%\citet{Brandstetter1906}	&	\citet{Kern1906}	&	\citet{Brandes1884}	\\
%\raggedright{\citet{Dempwolff1934,Dempwolff1937,Dempwolff1938}}	&	\citet{Jonker1914}	&	\citet{Friederici1913}	\\
%\raggedright{\citet{Salzner1960}}; Bima-Sumba 	&	\raggedright{\citet{KanskiKasprusch1931}}	&	\citet{Schmidt1926}	\\
%\citet{Haudricourt1965}	&		&	\citet{Cappell1943,Cappell1944}	\\
%\citet{Dyen1965,Dyen1978b}	&	\phantom{ }	&	\citet{Salzner1960}; Maluku	\\
%\citet{Dahl1976}	&		&	\citet{Blust1978,Blust1993,Blust2013}	\\
		%\lspbottomrule
	%\end{tabular}}
%\end{table}

%\begin{table}
	%\centering\caption{Previous studies grouping languages with the west or the east}
	%\label{tab:Previous studies grouping languages with the west or the east}
	%{\small{\wg\st{5pt}\begin{tabular}{p{0.3\columnwidth}p{0.3\columnwidth}p{0.3\columnwidth}}\lsptoprule
%Group with the west	&	Transitional	&	Group with the east	\\ \midrule
%\citet{Brandstetter1906}	&	\citet{Kern1906}	&	\citet{Brandes1884}	\\
%\citet{Dempwolff1934,Dempwolff1937,Dempwolff1938}	&	\citet{Jonker1914}	&	\citet{Friederici1913}	\\
%\citet{Salzner1960}; Bima-Sumba 	&	\citet{KanskiKasprusch1931}	&	\citet{Schmidt1926}	\\
%\citet{Haudricourt1965}	&		&	\citet{Cappell1943,Cappell1944}	\\
%\citet{Dyen1965,Dyen1978b}	&	\phantom{ }	&	\citet{Salzner1960}; Maluku	\\
%\citet{Dahl1976}	&		&	\citet{Blust1978,Blust1993,Blust2013}	\\
		%\lspbottomrule
	%\end{tabular}}
%\end{table}

Nearly 150 years ago the so-called “reverse genitive” was noted by
\citet{vanHoevell1877} for languages of Ambon-Lease in central Maluku.
We now know that it is typologically unusual
for SVO (subject-verb-object) languages with prepositions to have preposed possessors
(\citealp{Greenberg1966,Comrie1981}, \citealp[221]{Givon1984},
\citealp{Dryer2013a,Dryer2013b,Dryer2013d,Dryer2013f}).\footnote{\label{fn:01-03}
		In the WALS online sample of 488 SVO languages worldwide,
		only 43 have the typological combination of SVO (map 81A) + \it{prepositions}
		(map 85A) + \it{preposed possessors} (G-N; map 86A) \citep{Dryer2013b,Dryer2013d,Dryer2013f}.
		Of those 43 languages (< 9{\%}),
		13 are Austronesian languages within our target region and SHWNG.
		And if those 13 are correctly attributed to contact with typologically
		SOV Papuan languages that also have preposed possessors,
		then there remains only 30/488 other prepositional SVO languages with preposed possessors,
		or 6{\%} of the worldwide sample.
		The more common pattern for prepositional SVO languages is that of \it{postposed possessors} (N-G),
		for which WALS has 193/488 languages,
		or 40{\%} of the sample.
		The pattern of postposed possessors is the dominant pattern found
		in Austronesian languages west of what we are calling the Blust line in the typologically SVO, VSO, VOS,
		and languages with no dominant clause order of western Indonesia, Malaysia and the Philippines.
		In the worldwide sample,
		prepositional languages with postposed possessors combine to account for 323 languages,
		which is about ten times the number of prepositional SVO languages with preposed possessors.
		398/564 SOV languages in the worldwide sample have postposed possessors (71{\%}).}
Example \qf{01-01} shows the western order of \tsc{possessed + possessor} (post-posed possessor) from Malay.
Examples \qf{01-02} and \qf{01-03} show the eastern order of \tsc{possessor + possessed} (pre-posed possessor)
from Kotos Amarasi (west Timor)
and Buru (central Maluku) respectively.

\newpage
\begin{exe}
	\ex{western pattern: Malay}\label{ex:01-01}
	\sn{\gll	{} kəpala {} {} Yohanes {}\\
						[ head \bracr{\sub{POSSESSED}} [ John \bracr{\sub{POSSESSOR}}\\
			\glt	`John's head'}
	\ex{eastern pattern: Kotos Amarasi; \tsc{Timor-Babar}}\label{ex:01-02}
	\sn{\gll	{} naiʔ Yohanis {} ina {} ʔnaka-n {}\\
						[ Mr. John \bracr{\sub{POSSESSOR}} \tsc{3sg.nom} [ head-\tsc{3sg.gen} \bracr{\sub{POSSESSED}}\\
			\glt	`John's head'}
	\ex{eastern pattern: Buru; \tsc{Sula-Buru}}\label{ex:01-03}
	\sn{\gll	{} Yohanis {} {} olo=n {}\\
						[ John \bracr{\sub{POSSESSOR}} [ head=\tsc{3sg.gen} \bracr{\sub{POSSESSED}}\\
			\glt	`John's head'}
\end{exe}

Yet a number of languages in the region also retain the “western”
order of post-posed possessor, as in example \qf{01-04}.
These exceptions are discussed in more detail in \srf{02-03-01-01}.

\begin{exe}
	\ex{western pattern: Dhao and Havu; \tsc{Bima-Lembata}}\label{ex:01-04}
	\sn{\gll	{} kətu {} {} {\j}ohani {}\\
						[ head \bracr{\sub{POSSESSED}} [ John \bracr{\sub{POSSESSOR}}\\
			\glt	`John's head'}
\end{exe}

Early scholars simply noted the contrast in word order between languages
in eastern parts of the archipelago having the preposed possessor in contrast
with those to the west and north.
\citet{Brandes1884} used word order in possessive constructions
as a basis for dividing Austronesian languages into west
and east (with a line now known as the “Brandes line”),
thus grouping most of the languages of this region with those further east.
\citet{Friederici1913} supported the classificatory relevance of the preposed possessor,
and consequently also grouped the languages of this region with those to the east.
\citet{Jonker1914} rejected \citeauthor{Brandes1884}’
and \citeauthor{Friederici1913}’s use of word order as a basis for subgrouping,
and also considered other features that suggested Brandes’
line should be further west (\citeyear[367]{Jonker1914}).
\citet{Kern1906} and \citet{KanskiKasprusch1931} suggested the languages of
eastern Indonesia might be transitional between east and west.
\citet{Cowan1951} and \citet{Cappell1976}
continued to refer to the relevance of the preposed possessor.

\citet[235f]{Dyen1978b} provides a concise overview of
some early debates, and is worth quoting at length:

\begin{quote}
	``Three different propositions have been offered concerning these languages.
	The first is embodied in the Brandes Line which would divide these languages
	into an eastern and western group.
	According to \citeauthor{Brandes1884} the western boundary
	of the eastern group ran west of Roti in the south,
	and in the north, just west of the Sula Islands
	and thus east of Celebes and the Banggai Islands,
	and the Talaut and Sangir Islands.
	The basis of this division was the feature called ‘the preposed genitive’ which appears in languages east of this line.
	As an example \citeauthor{Brandes1884} cites Kisar \it{manu-keër} \it{‘bird egg’}
	(attribute plus head) as compared with Malay \it{telor ayam}
	\it{‘(egg chicken) chicken egg’} (head plus attribute).

	\citeauthor{KanskiKasprusch1931} (\citeyear{KanskiKasprusch1931}; hereafter \it{passim} KK) have reviewed
	the history of the discussion of the position of the languages of eastern Indonesia.
	After presenting \citeauthor{Brandes1884}’ view,
	they point out (\citeyear[883]{KanskiKasprusch1931})
	that in terms of \citeauthor{Brandes1884}’ criterion N.
	Adriani proposed that Banggai off Sulawesi should be included in this group.
	They themselves point out that Solorese on Flores must also be included.
	Furthermore Paulus Mitang, a Sikanese, has told me that Sika,
	or at least some dialects of Sika, must also be included.
	KK (\citeyear[884]{KanskiKasprusch1931}) go on to recall that \citeauthor{Schmidt1926}
	interpreted the feature of the preposed
	genitive as a relic of a Papuan substratum with which mixture had occurred
	and they attribute to Jonker (citing \citeauthor{Jonker1914} [\citeyear[4]{Jonker1914}])
	the belief that \citeauthor{Schmidt1926}
	assigned the languages of \citeauthor{Brandes1884}’ eastern group to the Melanesian languages.

	\citet{Schmidt1926} maintains \citeauthor{Brandes1884}’ eastern group,
	but has Sika at its western end and the Tanimbar
	and Aru languages at its eastern end.
	\citeauthor{Schmidt1926} does not mention the Austronesian languages of south Halmahera,
	but presumably they are to be included here.
	He does however include the New Guinea coastal languages among
	the so-called Papua-mixed group which he places among the Melanesian languages.
	Perhaps the fact that these coastal languages also show a
	‘preposed genitive’ led \citeauthor{Jonker1914} to conclude that \citeauthor{Schmidt1926}
	was extending the Melanesian boundary to include the eastern Indonesian languages.

	However this may be,
	Jonker for his part (\citeyear[263]{Jonker1914}) recommends ignoring the division
	of the languages of Indonesia into an eastern and western group.
	He sees no basis for either division.
	\citet[884]{KanskiKasprusch1931} on the other hand see the \citeauthor{Brandes1884}
	division justified on the basis that the eastern languages were influenced by the Papuan languages,
	not because of a radical linguistic difference from the western group
	nor because of an internal linguistic relationship within the eastern group.
	They claim that the eastern group constitutes a transitional group from the Indonesian languages to the Melanesian languages.
	They base their claim on an examination of the relations between the pronominal possessives in the languages.
	Their reasons for doing so is that these pronominal possessives together
	with the numerous further grammatical relations involved with them,
	are one of the strongest criteria for the classification of the
	Oceanic languages according to the best scholars of these languages.
	The statement of this claim is followed by the name Dempwolff in parentheses (\citeyear[884]{KanskiKasprusch1931})
	and he is presumably either among the scholars referred to or is the source of the opinion.

	The three views that have thus far been proposed are then as follows:
	
	\begin{itemize}
		\item[1)]{The original \citeauthor{Brandes1884} grouping which
							treats the eastern Indonesian languages as an Indonesian subgroup.
							The various additions to the original \citeauthor{Brandes1884} grouping are simple subvarieties of the same view.}
		\item[2)]{\citeauthor{Jonker1914}’s view that there is no eastern subdivision.
							This is perhaps best interpreted as claiming that
							there is neither a western subgroup nor an eastern subgroup.
							The implication appears to be that we simply have a large set of coordinate languages.} 
		\item[3)]{The view of \citeauthor{KanskiKasprusch1931} that the eastern languages of Indonesia
							form a subgroup which is different from the western subgroup
							and is transitional to the Melanesian or Oceanic languages.}'' \citep[235f]{Dyen1978b}
	\end{itemize}

\end{quote}

The general consensus today is that the preposed possessor in Austronesian languages
in the target region is the result of contact-induced change with typologically SOV
Papuan languages extant in the region before the arrival of the Austronesians,
and co-existing with them in intense contact over perhaps four millennia,
including bilingualism and intermarriage.
As a typological feature resulting from contact,
it has no bearing on genealogical relationships
and should not be used as a basis for linguistic subgrouping
(\citealp[290--292]{Grimes1991c}, \citealp{Grimes2000a,GrimesCIP,DonohueGrimes2008,
KlamerReesinkVanStaden2008,Schapper2015,Klamer2019,Schapper2024})

The preposed possessor is just one feature which shows an affinity between
the Austronesian languages of our target region with those of Melanesia.
There is an increasing awareness that many features that set these languages apart from those to the west
and north are due to intense contact with lexically
and typologically different SOV Papuan languages still found in the region.
And a growing awareness that prior to the arrival of Austronesian-speaking peoples,
there is trace evidence that Papuan languages were much more widespread in the region than they are today
(see \srf{01-06} and \srf{02-03-01}).\footnote{
		By Papuan we mean non-Austronesian language families,
		such as \tsc{Timor-Alor-Pantar} (associated with the \tsc{Trans New Guinea} languages),
		and \tsc{North Halmahera} languages.
		There were very likely other Papuan languages or language families
		in the region that have since disappeared \citep{Donohue2007}.}

The typological differences distinguishing the languages of Wallacea
from those further west is one of the hallmarks that led
\citet{Himmelmann2005} to divide Austronesian languages into
two major typological (not genealogical) subtypes:
the \it{symmetrical voice} type in the west and north,
and the \it{preposed possessor} type in the east,
each named after one typologically prominent feature.
This tells us that the languages in this region are different from those to the west
and north, and group together typologically,
but it does not address their genealogical affiliation.

\citet{Dyen1965,Haudricourt1965} and \citet{Dahl1976}
seem to include languages of Maluku (eastern Indonesia)
with the languages of the west rather than with those of Melanesia.
\citet{Dyen1978a} looked at data from Kei
and Kamarian in a wider discussion of languages in eastern Indonesia
and grouped them with the west.
\citet{Dyen1978b} looks at implications of vowels of
central Maluku languages for \tsc{PAn} subgrouping.

\it{Central Malayo-Polynesian} (CMP) as a purported subgroup within Austronesian is based on
\citet{Blust1974,Blust1978,Blust1979,Blust1981,Blust1982,
Blust1993,Blust2008,Blust2009,Blust2012b,Blust2013}.
It is roughly equivalent to \citeauthor{Esser1938}'s (\citeyear{Esser1938}) Ambon-Timor
combined with Bima-Sumba groups (followed by \citealp{Salzner1960}).
\citet{Dyen1965} proposed a “Moluccan Linkage”
(which was said to include “Sumba,
Sawu, Sikkic (including Endeh),
Fordatic, (including Jamden) Letic,
Buru, Ambic, Bonfia, Soboyo, Kuiwai”)
on the basis of his lexicostatistic calculations encompassing
approximately the same scope as Blust’s subsequent “CMP”.

\begin{figure}
	\centering
	\includegraphics[width=5cm]{AN-trad-tree.pdf}
	\caption{Austronesian tree adapted and updated from
					\citet{Blust1978,Blust1982,Blust1993,Blust1999,Blust2009}}\label{fig:01-02}
\end{figure}

The branching of Malayo-Polynesian in \frf{01-02} is referred
to by \citet[37]{Blust2009} as the “standard theory”,
which is a fair reflection of its acceptance as the status quo at that time.
With regard to an updated view of the Austronesian tree in \frf{01-02},
we note that “Western Malayo-Polynesian” (WMP) is no longer widely accepted as a cohesive subgroup \citep{Smith2017a}.
And “Central Malayo-Polynesian” (CMP) is also no longer widely accepted as a subgroup
\citep{Grimes1991a,PawleyRoss1993,Ross1995,Ross2008,Ross2017,
DonohueGrimes2008,GrayGreenhill2009,Schapper2015,Klamer2019,Zobel2024}

With regard to the east or west position of the languages in our target region,
Blust clearly groups “CMP” languages with the east under Central-Eastern Malayo-Polynesian (CEMP).
\citet[84]{Blust2013} sees a major division dividing the island of Sumbawa,
although the basis for that seems to be a subgrouping hypothesis,
without linguistic features actually specified,
“The presentation chosen here separates Sumbawa into western
and eastern halves,
since the major linguistic division separating
WMP from CMP languages runs through this island.”

\citet[note 13]{Blust1982} acknowledges,
“The proposed Central Malayo-Polynesian subgroup has thus far been
\it{assumed} without qualitative demonstration.
Dyen (1978) [\citeyear{Dyen1978a}] makes a similar assumption,
presumably on the basis of his \citeyear{Dyen1965}
lexicostatistical results.” [Emphasis ours].
The basis of \citet{Blust1982,Blust2012b} is that languages in the “CMP”
region share with EMP languages reflexes of two
lexical innovations when first encountering marsupials east of the Wallace line
(i.e. CEMP *kandoRa ‘cuscus’ (\frf{01-03}),
and *mans(aə)r ‘bandicoot’,
(\frf{01-04}); note the two pictures are not equal in scale).

%%%%%%%%%%%%%%%%%%%%%%%%%%%%%%%%%%%%%%%%%%%%%%%%%%%%%%%%%%%%%%%%%%%%%%%&%
%%                                                                     %%
%% NOTE FOR TYPESETTING: I have added some manual spacing to have      %%
%% the pictures not too close together, but also not too far apart     %%
%% and to have them appear on the same page as the text referencing    %%
%% them. The sizes have also been reduced (they were too big in        %%
%% the proofreading version.                                           %%
%%                                                                     %%
%% Please consult before editing how these pictures have been typeset! %%
%% Changing the sizes could have bad flow on effects for pagination.   %%
%%                                                                     %%
%%%%%%%%%%%%%%%%%%%%%%%%%%%%%%%%%%%%%%%%%%%%%%%%%%%%%%%%%%%%%%%%%%%%%%%%%

\begin{figure}\vspace{-2mm}\centering
  \subfigure[Gray cuscus \newline (\it{Phalanger orientalis})]{\label{fig:01-03}
    \includegraphics[width=.35\textwidth]{01-03.jpg}
    }\hspace{1cm}
    \subfigure[Bandicoot \newline (\it{Peramelidae sp.})]{\label{fig:01-04}
    \includegraphics[width=.25\textwidth]{01-04.jpg}
    }
    \caption{Marsupials}\vspace{-3mm}
\end{figure}

%\begin{figure}
  %\subfigure[Gray cuscus (\it{Phalanger orientalis})]{\label{fig:01-03}
    %\includegraphics[width=.45\textwidth]{01-03.jpg}
    %}
    %\subfigure[Bandicoot (\it{Peramelidae sp.})]{\label{fig:01-04}
    %\includegraphics[width=.45\textwidth]{01-04.jpg}
    %}
    %\caption{\missingcaption}
%\end{figure}

These etyma were not found in WMP languages,
thus implying a CEMP subgroup that encompasses both EMP
and the languages in the CMP region.
The assumptions behind the first article,
along with the interpretation of the data
and its implications for subgrouping have been challenged by \citet{Schapper2011}.
We return briefly to methodological issues that
have complicated this discussion in \srf{02-04},
and at length in \crf{ch:Mar}.

Although \citet{Stresemann1927,Chlenov1976,Collins1981,Collins1983,
Collins1986,Hughes1987,LoskiLoski1989,Mills1991,Taber1993}
and others have looked at subgrouping within limited areas of Maluku
(with varying methodologies and often conflicting conclusions),
nobody has really looked at the over-all relationships of the languages of Nusa Tenggara
and Timor with the languages of Maluku with a view to internal subgrouping
by looking systematically at the historical phonology in a principled way.
Thus, we have had no credible bird’s-eye view of first-order divisions within
“CMP” that withstands the scrutiny of additional data or careful methodology.

\citet{Blust1993} brings the higher level notion of “CMP”
out of the realm of educated impressions
and into the realm of scientific discussion,
and is a substantial contribution to our understanding
of the \tsc{An} languages of eastern Indonesia.
But in spite of his declared intent
“to argue for the reality of a Central Malayo-Polynesian subgroup”
(\citeyear[244]{Blust1993}),
and vigorous defences of such a subgroup,
giving impressions that have misled several to assume that CMP
is a “branch” of PMP whose daughter languages reflect
the characteristics of a parent Proto-CMP,\footnote{
		As evidenced in several emails
		and conversations with established
		and emerging Austronesian scholars over the past few years.}
nevertheless a careful reading of \citet{Blust1993} supports only
the weaker claim that there are non-overlapping innovations found in the region,
some of which are not found elsewhere.
A careful examination of the fuller data reveals that many of these are sporadic,
not replacement innovations,
and are also found outside the “CMP” region,
including to its west.

We use the label “Blust line”
to refer to the boundary which demarcates Blust's
putative CMP and CEMP subgroups
from other Austronesian languages.
The Blust line is the boundary
between the languages to the west
and north of our Wallacean subgroups (i.e.
outside our target region).

\citet[32]{Blust2013} explains the inability to find exclusively
shared patterned innovations within the region as follows: “Many of the lexical
and phonological innovations used to assign languages
to CMP do not cover the entire set of languages,
and so suggest that this grouping arose from an original dialect chain
that served as a ‘diffusion corridor’ rather than through a series of ‘clean’ language splits.”
In spite of the significant contributions of
\citet{Blust1993,Blust2008,Blust2009,Blust2012b,Blust2013},
we are still left with important questions that remain unanswered:

\begin{itemize}
	\item[1.] What defines a “CMP” language? What diagnostic features should we be looking for?
	\item[2.] What distinguishes a “CMP” language from non-CMP languages?
						How do we know what is not a “CMP” language?
	\item[3.] What subgroups are found within the purported “CMP” subgroup?
	\item[4.] How are these purported “CMP” subgroups related to each other?
						How do they link up within “CMP”, or to a common parent language or linkage?
\end{itemize}

The strongest arguments -- considered best practice in comparative linguistics -- are
when one can demonstrate \it{exclusively shared innovations} in the historical phonology,
the morphosyntax, and the lexicon.
An innovation occurs in patterned ways within the target group,
and does not occur in those same patterned ways outside the target group.
Lexical items are easily borrowed,
so lexical data are not given as much weight as patterns in the historical phonology.
And patterns in the historical phonology must be established
before one can argue for shared innovations in the morphosyntax.
If the feature occurs in patterned ways outside the target group of languages
--- particularly nearby ---
that feature cannot be said to be an exclusively shared innovation.
If the feature occurs only sporadically within the target group of languages
and not in patterned ways,
that feature cannot be claimed to be an exclusively shared innovation
deriving from a common parent language.
Sporadic patterns that align with other patterns in their distribution
can form a bundle that contributes towards a more convincing picture,
but of themselves do not constitute subgrouping arguments.
Sporadic patterns that do not align with other patterns are simply sporadic patterns,
and can easily be due to parallel development,
or contact.

\citet{Blust1993,Blust2013} maintains a weaker view,
that his purported \it{Central Malayo-Polynesian} (CMP) subgroup is based on
“overlapping distributions of innovated features”
developing from a linkage (\citeyear[242]{Blust1993}).
A linkage arises,
not from a single shared parent language,
but “as languages progressively differentiate out of a dialect network” \citep{Ross2017}.
This view is still not an adequate explanation by itself for the diverse data
and sporadic patterns we find throughout the so-called “CMP” region,
as we show in the following chapters.

Keeping in mind that “the classical comparative method
does not readily recognise [{\ldots] contact-induced change” \citep[49]{Ross1995},
all of this in a region of intense contact with Papuan languages over
four millennia helps to explain Blust’s frustrations with regard to
“CMP” for which \citeauthor{Blust2013} acknowledges (\citeyear[32]{Blust2013}),
“Higher-order subgrouping within CMP has proven difficult to date.”
and “Many of the lexical
and phonological innovations used to assign languages to
CMP do not cover the entire set of languages,…”.\footnote{
		\citet{GrayGreenhill2009},
		looking at phylogenetics from the data in 
		\citet{GreenhillGray2008} conclude that CMP is “not supported”.
		An implication of rethinking the legitimacy of
		Western Malayo-Polynesian (WMP) \citep{Smith2017b}
		and CMP, is that if WMP
		and CMP are not valid subgroups,
		then many etyma that have been ascribed at the level of PMP,
		but whose reflexes are not found in the Philippines or much of western Indonesia,
		need to be re-examined
		and reassigned at a different level.
		Note that \citet[754--756]{Blust2013}
		does not include CMP in his Austronesian tree.}
Regarding the higher-level CEMP,
\citet[75]{Blust2009} also acknowledges that it “poses some of the most complex
and challenging subgrouping issues that are found in [Austronesian].”

No serious scholar currently defends the stronger view that the
Austronesian languages of eastern Indonesia
and Timor-Leste (excluding the \tsc{South Halmahera-West New Guinea} languages)
form a classic subgroup with a single parent language in which the daughter
languages manifest a number of exclusively
shared innovations in their historical phonologies,
their morphosyntax,
and their lexicons that identify languages within the group,
and distinguish them from languages outside the group.
\citet[92]{Blust2008} admits,
“The historical phonology of the languages of
the Lesser Sunda islands seems almost calculated to perplex.”
The data simply do not support the stronger view of a classic subgroup.

Since various things that have been proposed do not meet standard criteria for subgrouping arguments,
for the moment what we are left with is mostly a geographic region,
and a negative definition: \it{“CMP” languages are “CEMP” languages that are not EMP languages}.
This is not very satisfying scientifically.
\citet{Grimes1991a,PawleyRoss1993,Ross1995,Ross2008,
Ross2017,GrayGreenhill2009,DonohueGrimes2008,Schapper2015,Klamer2019,Zobel2024}
conclude that the notion of “CMP” as a subgroup
within a family tree model is not well supported.

Arguments published in the 1970s--1990s proposing localised
subgroups all require a necessary refinement.
Due to recent advances in the field,
as well as access to better data (in both quality and quantity).
Some of these earlier proposals were based on outdated mergers
of PMP *D/*d, *Z/*z, *-ay/*-əy.\footnote{
		The distinctions between *Z
		and *z, and *-ay and *-əy have been abandoned by Austronesian scholars.
		\citet[582--584]{Blust2013} only maintains the distinction between *D
		and *d word finally for a few languages in western Indonesia,
		noting that he does so uneasily.}
Some did not look at geographically near neighbours across political boundaries.
In addition there was often the assumption of a monolithic Austronesian dispersal.
All of these issues are looked at in subsequent chapters.

There remain significant as yet unanswered questions
regarding the macro- and micro-classification of the languages in this region.
It is hoped that the data
and arguments presented in this study will contribute positively to the discussion,
and to a credible refinement of the linguistic picture of eastern Indonesia
and Timor-Leste.

Discussion of the history of the study of individual subgroups
in the region are found in the relevant chapters.
Here we simply note that people have not been able to credibly
link anyone’s purported subgroups in the region together
on the basis of the historical phonology,
and many proposed lexical “innovations” are problematic,
are not found in all subgroups,
are not replacement innovations, are not exclusive,
and do not line up with other proposed lexical innovations.
A convincing bundle of exclusively shared innovations
which tell a coherent story have not been proposed.

Before moving on,
we want to pay tribute to Erwin Stresemann,
a German ornithologist who published a grammar of Paulohi
on south-central Seram \citep{Stresemann1918} with additional
comparative notes on Seram and Buru.
He later returned to Germany
and worked side-by-side with J. C. G. Jonker
and Otto Dempwolff where he produced a volume
\citet{Stresemann1927} whose significance is perhaps often overlooked
in the history of comparative Austronesian studies.

\begin{itemize}
	\item \citet{Stresemann1927} predates the works of
				\citet{Dempwolff1934,Dempwolff1937,Dempwolff1938} that get much more attention.
	\item \citeauthor{Stresemann1927}'s (\citeyear{Stresemann1927})
				reconstructions for Central Maluku, which are
				accessible as “PAMS” in \citet{WurmWilson1983},
				form the basis for many of \citeauthor{Collins1983}' (\citeyear{Collins1983})
				reconstructions for Proto-Central Maluku (PCM),
				and for some of Blust’s PCMP reconstructions \citet{BlustTrussel2020}.
				With more and better data from more languages,
				we (collectively) are now able to add some final consonants here
				and there, but the building blocks are present in \citeauthor{Stresemann1927}’s work,
				e.g. \citeauthor{Stresemann1927}’s distinction between PCM **v
				and **b tracks the difference between PMP *b
				and *mb respectively,
				a distinction which must be preserved throughout the target region.
	\item Some of \citet{Stresemann1927}'s reconstructions for PMP
				remain unchanged to this day.
				For example if we update the orthography to current usage we have PMP *bulaw-an ‘gold’,
				*pajay ‘rice’, *pulut ‘glue, paste, adhesive’,
				*səŋət ‘mosquito, stinging fly’, *suja ‘mantrap’, and *təbus ‘ransom, redeem’, among others.
				In some cases, the gloss has since been adjusted.
	\item \citet[4]{Collins1983} notes that \citeauthor{Stresemann1927}'s (\citeyear{Stresemann1927})
				effort to define “Ur-Ambon” as a distinct subgroup,
				along with lower level “Sub-Buru”, “Sub-Ambon”, and “Sub-Seran”,
				by identifying splits and mergers in the historical phonology
				and some morphology,
				was “one of the earliest systematic attempts at subgrouping in the [Austronesian] family.”
				While \citet{Collins1983} pointed out a number of exceptions
				to \citeauthor{Stresemann1927}'s (\citeyear{Stresemann1927}) proposed mergers requiring
				some of the distinctions to be preserved in the historical phonology of these languages,
				many of \citeauthor{Stresemann1927}’s observations nevertheless accurately
				summarise the patterns for many languages in the region.
				Given what little he had to work with,
				it is quite impressive.
	\item \citet[90]{Blust2013} puts it this way,
				“In many ways this book [\cite{Stresemann1927}] still
				stands as the single most important contribution
				to the comparative study of these languages,
				although important additional studies,
				particularly in the area of subgrouping,
				have been done more recently by James T.
				Collins \citep{Collins1982d,Collins1983}.”
\end{itemize}

We also note that \citet{Jonker1906} is an early attempt at comparative morphosyntax.
The analysis of the suffixes and enclitics that he investigated continue
to challenge modern investigators
(\citealp[vi]{Fox1968}, \citealp[30--33]{Collins1983}, \citealp{Grimes1991b}).}
